\begin{abstract}
\addchaptertocentry{\abstractname} % Add the abstract to the table of contents
\bigskip
\noindent This thesis investigates tendon-driven mechanisms for robotic manipulation, focusing on four mechanism-level challenges: maintaining stable tension transmission, achieving dexterous motion with low distal mass, preserving compliant and adaptive contact, and operating reliably under geometric constraints. Rather than presenting grasping and surgical instrumentation as separate topics, the thesis develops a transferable tendon-driven design framework and validates it through two embodiments.

The first embodiment is Lasso Gripper, a novel string-driven end-effector inspired by the traditional lasso and the uurga. Instead of relying on concentrated antipodal contact, the mechanism launches and retracts a controllable string loop to form an adaptive capture region around the target. This architecture expands the effective capture range, distributes contact force along the tensile element, and improves the handling of oversized, irregularly shaped, and fragile objects. Lasso Gripper integrates a dedicated launch-and-retraction mechanism, coordinated motor actuation, grasp planning, and real-time tension regulation. Experimental validation demonstrates successful grasping of diverse static and dynamic targets, including animal figures, vegetables, balloons, and moving objects, while also showing clear advantages over conventional rigid-finger grasping in shape adaptability and gentle contact.

Beyond demonstrating a new gripper, Lasso Gripper serves as a mechanism study platform from which transferable design principles are extracted. These include proximal actuation for reducing distal inertia, differential tendon routing for compact multi-degree-of-freedom behavior, and explicit tension management for improving repeatability and force transmission. The thesis then translates these principles into a second embodiment: a handheld motorized laparoscopic instrument for minimally invasive surgery.

The surgical instrument validates the same tendon-driven framework under far stricter geometric and ergonomic constraints. It employs a compact arrangement of input sensing, proximal actuation, tendon routing, and tension-preserving reel architecture to realize multiple distal degrees of freedom in a hand-held form factor. Its control stack combines signal filtering with closed-loop motor feedback to suppress disturbance and improve motion fidelity. In this way, the thesis shows that the mechanism principles first clarified through adaptive robotic grasping can be transferred to clinically constrained surgical manipulation.

Taken together, the two embodiments establish the central contribution of this thesis: tendon-driven mechanisms can be systematically designed as a transferable framework for adaptive grasping, dexterous manipulation, and safe interaction across both general robotic and surgical domains.

\begin{flushright}
   (332 words)
\end{flushright}


\end{abstract}
